\documentclass[printcopy,fontset=windows]{../Style/neuthesis}
\usepackage[bibtex,myhdr,table,list,geometry,math]{../Style/artratex}
\usepackage{booktabs}
\usepackage{tabularx}
\usepackage{threeparttable}
\usepackage{array}
\usepackage{etoolbox}

\AtBeginEnvironment{table}{\footnotesize}
\AtBeginEnvironment{tabular}{\footnotesize}
\AtBeginEnvironment{tabularx}{\footnotesize}
\AtBeginEnvironment{tablenotes}{\footnotesize}

\renewcommand{\arraystretch}{1.18}
\setlength{\tabcolsep}{4pt}
\newcolumntype{Z}{>{\raggedright\arraybackslash}X}

\begin{document}

\begin{table}[p]
  \centering
  \bicaption{HSSTT 章节对比方法总览}{Overview of comparison methods in the HSSTT chapter}
  \label{tab:hsstt_methods}
  \begin{threeparttable}
    \begin{tabularx}{\textwidth}{>{\raggedright\arraybackslash}p{0.18\textwidth}Z}
      \toprule
      方法 & 方法实现说明 \\
      \midrule
      MRVR &
      通过稀疏贝叶斯回归拟合输入到 SOH 的映射，并用参数后验给出预测分布 \\
      GP-ICE &
      通过高斯过程对输出函数建模，并借助核函数先验刻画预测不确定性 \\
      BMNN &
      通过 Bayes by Backprop 学习权重后验，再由参数采样形成预测分布 \\
      SACN &
      通过胶囊网络聚合时空特征，并借由特征响应的变化间接表征不确定性 \\
      DeNN &
      通过多个子模型的集成或融合统计预测分歧，以此近似不确定性 \\
      BDAE &
      通过随机失活下的多次前向传播形成预测集合，并用输出波动量化不确定性 \\
      SNGP &
      通过谱归一化稳定特征映射，并在输出端接高斯过程头近似后验分布 \\
      \bottomrule
    \end{tabularx}
  \end{threeparttable}
\end{table}
\clearpage

\begin{table}[p]
  \centering
  \bicaption{RODTNet 章节对比方法总览}{Overview of comparison methods in the RODTNet chapter}
  \label{tab:rodtnet_methods}
  \begin{threeparttable}
    \scriptsize
    \begin{tabularx}{\textwidth}{>{\raggedright\arraybackslash}p{0.18\textwidth}Z}
      \toprule
      方法 & 方法实现说明 \\
      \midrule
      iTransformer &
      通过将变量维度作为建模对象，把多变量序列映射到反向注意力结构中 \\
      PatchTST &
      通过将长序列切分为 patch 并送入 Transformer，以局部片段建模全局依赖 \\
      TSMixer &
      通过纯 MLP 在时间维与通道维上交替混合特征，不依赖显式注意力 \\
      PINN &
      通过把物理约束写入损失函数，引导网络学习满足机理先验的退化映射 \\
      MambaLithium &
      通过状态空间选择性扫描机制建模长序列退化依赖 \\
      DeNN-bagging &
      通过自助采样训练多个子模型，并对输出结果做平均融合 \\
      DeNN-fusion &
      通过特征或预测级融合多个子模型，以提高跨样本稳定性 \\
      DeNN-voting &
      通过对子模型输出进行投票汇总，降低单模型偏置 \\
      LSTM-DA &
      通过 LSTM 提取时序特征，并配合域对齐模块缩小源域和目标域差异。 \\
      ConvMMD &
      通过卷积网络提取特征，并用 MMD 对齐不同域的分布。 \\
      DKTL &
      通过核映射增强特征可分性，并结合迁移学习实现跨域建模。 \\
      DDAN &
      通过域判别器对抗训练，使特征提取器学习域不变表示。 \\
      DR-Net &
      通过重构退化模式来增强特征表达，并抑制域间差异。 \\
      CRNN &
      通过卷积层提取局部模式，再用循环层建模时序依赖。 \\
      MAGNet &
      通过图结构传播退化关联，并在域间共享图表示。 \\
      iTransformer+ &
      先在源域预训练，再直接在目标域上微调 iTransformer \\
      PatchTST+ &
      先在源域预训练，再直接在目标域上微调 PatchTST \\
      TSMixer+ &
      先在源域预训练，再直接在目标域上微调 TSMixer \\
      PINN+ &
      先在源域预训练，再直接在目标域上微调 PINN \\
      \bottomrule
    \end{tabularx}
  \end{threeparttable}
\end{table}
\clearpage

\begin{table}[p]
  \centering
  \bicaption{DA-TMLLM 章节对比方法总览}{Overview of comparison methods in the DA-TMLLM chapter}
  \label{tab:da_tmllm_methods}
  \begin{threeparttable}
    \scriptsize
    \begin{tabularx}{\textwidth}{>{\raggedright\arraybackslash}p{0.18\textwidth}Z}
      \toprule
      方法 & 方法实现说明 \\
      \midrule
      VQVAE &
      通过离散码本对信号进行压缩重构，以生成更平滑的少数类样本 \\
      DCGAN &
      通过卷积生成器与判别器对抗训练，学习原始信号分布 \\
      QSCGAN &
      通过条件生成与判别联合约束，按类别条件生成目标信号 \\
      BearLLM &
      通过大语言模型把文本先验与信号表征结合起来生成诊断结果 \\
      MLMM &
      通过多模态表示对齐，将文本提示与振动信号联合建模 \\
      WDSN &
      通过显式检测分支区分已知类别输入与非预期输入 \\
      MDCC &
      通过校正模块修正检测偏差，并抑制非预期输入的误判 \\
      LLM-Check &
      通过语言模型生成的语义一致性来判断输入是否符合预期 \\
      MIND &
      通过多模态融合后的判别头区分已知类别与异常输入 \\
      SHalluDetect &
      通过结构化规则和检测器联合识别幻觉输入 \\
      LLM-OOD &
      通过大语言模型估计输入是否偏离训练分布 \\
      DA-TMLLM &
      通过关键 token 驱动的不确定性重加权，把文本推理与信号判别联合起来 \\
      \bottomrule
    \end{tabularx}
  \end{threeparttable}
\end{table}

\end{document}
